% !TEX root = ../main.tex
\begin{ouverture}[seconddegre]{1}{Polynômes du second degré}{Partie 1 · Algèbre}
{Formes d'un trinôme\enspace·\enspace Équations\enspace·\enspace Signe\enspace·\enspace Modélisation}
\objectif{Choisir entre forme développée, factorisée et canonique.}
\objectif{Résoudre une équation et une inéquation du second degré.}
\objectif{Lire une parabole et déterminer son sommet.}
\objectif{Modéliser un problème d'aire, de coût ou de trajectoire.}
\prerequis{Identités remarquables et équations produit nul.}
\prerequis{Intervalles et repérage dans le plan.}
\end{ouverture}
\essentiel
\essSection{1}{Trois écritures, trois usages}
\begin{definition}[Trinôme]
Une fonction du second degré s'écrit $f(x)=ax^2+bx+c$, avec $a\neq0$.
Sa courbe est une parabole. La forme développée donne immédiatement $f(0)=c$.
\end{definition}
\begin{propriete}[Formes utiles]
La forme canonique $a(x-\alpha)^2+\beta$ donne le sommet $S(\alpha\pv\beta)$
et l'axe de symétrie $x=\alpha$. La forme factorisée $a(x-r_1)(x-r_2)$, lorsqu'elle
existe, donne les racines $r_1$ et $r_2$.
Alors $r_1+r_2=-b/a$ et $r_1r_2=c/a$, comme on le vérifie en développant.
Les fonctions du second degré qui s'annulent en deux réels distincts $r_1$ et $r_2$
sont les fonctions $x\mapsto a(x-r_1)(x-r_2)$, avec $a\neq0$.
\end{propriete}
\begin{exemple}
$x^2-4x+3=(x-2)^2-1=(x-1)(x-3)$.
Les racines sont 1 et 3 ; le minimum vaut $-1$, atteint en 2.
Pour une forme développée $ax^2+bx+c$, le sommet a pour abscisse $\alpha=-b/(2a)$
et pour ordonnée $f(\alpha)$. On peut aussi compléter un carré dans les cas simples.
La courbe de $x\mapsto f(x-m)$ se déduit de celle de $f$ par une translation de $m$
unités parallèlement à l'axe des abscisses : la parabole de $(x-2)^2-1$ est celle de $x^2$
décalée de 2 vers la droite et de 1 vers le bas.
\end{exemple}
\begin{center}
\begin{tikzpicture}\begin{axis}[otLivreGraphe,width=104mm,height=63mm,xmin=-0.5,xmax=4.5,ymin=-2,ymax=5,xtick={0,1,2,3,4},ytick={-1,1,3,5}]
\addplot[otNavy,thick,domain=-0.4:4.4,samples=90]{(x-2)^2-1};
\addplot[otGreen,dashed] coordinates {(2,-2)(2,5)};
\addplot[only marks,otNavy] coordinates {(1,0)(3,0)(2,-1)};
\node[below right,font=\small] at (axis cs:2,-1){$S(2;-1)$};
\end{axis}\end{tikzpicture}\par
\legende{La parabole de $f(x)=x^2-4x+3$ et son axe de symétrie.}
\end{center}
\begin{attention}
Une fonction peut ne pas avoir de racine réelle. La forme factorisée sur $\R$
n'existe alors pas. La forme canonique existe toujours.
\end{attention}
\clearpage
\essSection{2}{Résoudre une équation}
\begin{theoreme}[Discriminant]
Pour $ax^2+bx+c=0$, on pose $\Delta=b^2-4ac$.
Si $\Delta>0$, les deux solutions sont $\dfrac{-b-\sqrt\Delta}{2a}$ et
$\dfrac{-b+\sqrt\Delta}{2a}$. Si $\Delta=0$, la seule solution est
$-\dfrac b{2a}$. Si $\Delta<0$, il n'y a aucune solution réelle.
\end{theoreme}
\begin{savoirfaire}[Choisir le calcul le plus court]
Avant le discriminant, chercher une factorisation évidente :
$x^2-9=(x-3)(x+3)$, ou $2x^2-6x=2x(x-3)$.
Si une racine $r_1$ est évidente, l'autre se déduit de $r_1r_2=c/a$.
Pour compléter un carré, utiliser $x^2+2ux=(x+u)^2-u^2$.
\end{savoirfaire}
\essSection{3}{Étudier un signe et résoudre une inéquation}
\begin{propriete}
Lorsque deux racines $r_1<r_2$ existent, $f$ est du signe de $a$ à l'extérieur
de $[r_1;r_2]$, du signe opposé entre les racines, et s'annule aux racines.
Si $\Delta<0$, $f$ garde le signe de $a$. Si $\Delta=0$, $f$ garde le signe
de $a$ sauf à sa racine où elle vaut 0.
\end{propriete}
\begin{center}\begin{tableau}
\otTabL[2.6]{2.7}{$x$/0.8,$(x-1)(x-3)$/0.9}{$-\infty$,$1$,$3$,$+\infty$}
\tkzTabLine{,+,z,-,z,+,}
\end{tableau}\end{center}
\begin{savoirfaire}[Passer du signe à l'ensemble des solutions]
$x^2-4x+3\leq0\iff x\in[1;3]$.
Pour une inégalité stricte, les racines sont exclues :
$x^2-4x+3<0\iff x\in]1;3[$.
\end{savoirfaire}
\begin{attention}
Dans un quotient, interdire les valeurs qui annulent le dénominateur,
même si une simplification paraît possible. Pour un problème concret,
intersecter les solutions avec le domaine du modèle.
\end{attention}
\exercices

\rubrique{Découverte}
\begin{blocLivre}{ex-1-1}
\begin{exercice}[1]{Reconnaître un trinôme}
\begin{questions}
\item Parmi $3x^2-2x+5$, $4x-1$, $(x-2)^2$, $x^3+x$, repérer les polynômes du second degré.
\item Pour chacun d'eux, donner les coefficients $a$, $b$ et $c$.
\end{questions}
\end{exercice}
\end{blocLivre}
\begin{blocLivre}{ex-1-2}
\begin{exercice}[1]{Développer ou factoriser}
\begin{questions}
\item Développer $(x-3)(x+2)$ et $2(x-1)^2-8$.
\item Factoriser $x^2-25$ et $3x^2-12x$.
\end{questions}
\end{exercice}
\end{blocLivre}
\begin{blocLivre}{ex-1-3}
\begin{exercice}[1]{Lire la forme canonique}
Soit $f(x)=2(x-3)^2-8$.
\begin{questions}
\item Donner le sommet, l'axe de symétrie et le minimum de $f$.
\item Calculer $f(0)$ et résoudre $f(x)=0$.
\end{questions}
\end{exercice}
\end{blocLivre}
\begin{blocLivre}{ex-1-4}
\begin{exercice}[1]{Calculer un discriminant}
Résoudre dans $\R$ :
\begin{tasks}(3)
\task $x^2-5x+6=0$.
\task $x^2+4x+4=0$.
\task $2x^2+x+2=0$.
\end{tasks}
\end{exercice}
\end{blocLivre}
\begin{blocLivre}{ex-1-5}
\begin{exercice}[1]{Signe d'une forme factorisée}
Soit $P(x)=-2(x+1)(x-4)$.
\begin{questions}
\item Dresser le tableau de signes de $P$.
\item Résoudre $P(x)\geq0$, puis $P(x)<0$.
\end{questions}
\end{exercice}
\end{blocLivre}
\clearpage
\rubrique{Entraînement}
\begin{blocLivre}{ex-1-6}
\begin{exercice}[2]{Choisir la bonne forme}
On pose $f(x)=x^2-6x+5$.
\begin{questions}
\item Montrer que $f(x)=(x-3)^2-4$ et que $f(x)=(x-1)(x-5)$.
\item Résoudre $f(x)=5$, puis $f(x)\leq0$.
\item Résoudre $f(x)\leq-3$ en choisissant une forme adaptée.
\end{questions}
\end{exercice}
\end{blocLivre}
\begin{blocLivre}{ex-1-7}
\begin{exercice}[2]{Deux paraboles}
On considère $f(x)=x^2-4x+3$ et $g(x)=-x^2+2x+3$.
\begin{center}\begin{tikzpicture}\begin{axis}[otLivreGraphe,width=101mm,height=48mm,xmin=-1,xmax=4,ymin=-2,ymax=6,xtick={-1,0,1,2,3,4},ytick={-2,2,4,6}]
\addplot[otNavy,thick,domain=-0.6:4,samples=75]{x^2-4*x+3};
\addplot[otGreen,thick,domain=-1:3.4,samples=75]{-x^2+2*x+3};
\node[otEtiquette,text=otNavy] at(axis cs:3.7,2.7){$\mathcal C_f$};
\node[otEtiquette,text=otGreenDark] at(axis cs:2,4.7){$\mathcal C_g$};
\end{axis}\end{tikzpicture}\end{center}
\begin{questions}
\item Conjecturer les abscisses des points d'intersection.
\item Les calculer exactement, puis déterminer quand $f(x)\leq g(x)$.
\end{questions}
\end{exercice}
\end{blocLivre}
\begin{blocLivre}{ex-1-8}
\begin{exercice}[2]{Un quotient et une valeur interdite}
Résoudre $\dfrac{x^2-5x+6}{x-1}\geq0$.
Préciser le domaine, factoriser le numérateur, dresser un tableau de signes
et écrire l'ensemble des solutions avec les bons crochets.
\end{exercice}
\end{blocLivre}
\begin{blocLivre}{ex-1-9}
\begin{exercice}[2]{Retrouver une fonction}
Une parabole coupe l'axe des abscisses en $-2$ et $4$, et passe par $A(0;-8)$.
\begin{questions}
\item Justifier que sa fonction s'écrit $f(x)=a(x+2)(x-4)$, puis calculer $a$.
\item Développer $f$ et déterminer son sommet.
\item Résoudre $f(x)=-5$.
\end{questions}
\label{figlivre-r01-09-parabole-ex}
\figureLivre{r01-09-parabole}

\end{exercice}
\end{blocLivre}
\begin{blocLivre}{ex-1-10}
\begin{exercice}[2]{Somme et produit des racines}
\begin{questions}
\item Deux réels ont pour somme 7 et pour produit 10. Construire une équation
dont ils sont les racines et retrouver ces nombres.
\item Même question avec une somme de 4 et un produit de 5. Que peut-on conclure ?
\item Factoriser mentalement $x^2+x-12$ en cherchant les deux racines.
\end{questions}
\end{exercice}
\end{blocLivre}
\begin{blocLivre}{ex-1-11}
\begin{exercice}[2]{Discuter selon un paramètre}
Pour $m\in\R$, on considère $x^2-4x+m=0$.
\begin{questions}
\item Écrire $x^2-4x+m$ sous forme canonique.
\item Selon $m$, déterminer le nombre de solutions réelles.
\item Donner les solutions pour $m=3$, $m=4$ et $m=5$.
\item Pour quelles valeurs de $m$ l'expression est-elle positive ou nulle pour tout $x$ ?
\end{questions}
\end{exercice}
\end{blocLivre}
\clearpage
\rubrique{Problèmes}
\begin{blocLivre}{ex-1-12}
\begin{exercice}[3]{Un jardin contre un mur}
On dispose de 40 m de clôture pour fermer trois côtés d'un rectangle,
le quatrième étant un mur. On note $x$ la largeur perpendiculaire au mur.
\begin{center}\begin{tikzpicture}[x=1cm,y=1cm]
\fill[otNavyTint](0,0)rectangle(6,1.7);
\draw[otNavy,thick](0,1.7)--(0,0)--node[below]{$40-2x$}(6,0)--node[right]{$x$}(6,1.7);
\draw[otMuted,line width=2pt](-0.4,1.7)--(6.4,1.7);
\node[above,font=\small]at(3,1.7){mur};
\end{tikzpicture}\par\legende{Schéma de la situation, non à l'échelle.}\end{center}
\begin{questions}
\item Justifier $0<x<20$ et montrer que l'aire vaut $A(x)=-2x^2+40x$.
\item Écrire $A$ sous forme canonique ; en déduire l'aire maximale et les dimensions du jardin.
\item Pour quelles largeurs l'aire est-elle d'au moins $150$ m$^2$ ?
\end{questions}
\end{exercice}
\end{blocLivre}
\begin{blocLivre}{ex-1-13}
\begin{exercice}[3]{Une vente rentable}
Pour $q\in[0;100]$, on modélise le prix de vente unitaire par $p(q)=50-0{,}2q$
et le coût total par $C(q)=10q+500$, en euros.
\begin{questions}
\item Exprimer la recette $R(q)$ et le bénéfice $B(q)$.
\item Vérifier $B(q)=-0{,}2(q-100)^2+1500$ ; déterminer le bénéfice maximal sur le domaine.
\item Résoudre $B(q)\geq0$ et en déduire les quantités entières rentables. Donner une approximation à $10^{-2}$ des seuils utiles.
\end{questions}
\end{exercice}
\end{blocLivre}
\begin{blocLivre}{ex-1-14}
\begin{exercice}[3]{Un ballon au-dessus d'un filet}
La hauteur d'un ballon, en mètres, est $h(t)=-5t^2+10t+1$,
où $t\geq0$ désigne le temps en secondes après le lancer.
Le modèle est utilisé jusqu'au premier contact avec le sol.
\begin{questions}
\item Déterminer la hauteur initiale et la hauteur maximale.
\item Résoudre $h(t)=0$ ; à quel instant le ballon touche-t-il le sol ?
\item Déterminer les instants où $h(t)\geq4$.
\item Le ballon atteint le plan vertical d'un filet de 4 m à $t=1{,}5$ s.
Le franchit-il ? Justifier par un calcul et interpréter le résultat.
\end{questions}
\label{figlivre-r01-14-ballon-ex}
\figureLivre{r01-14-ballon}

\end{exercice}
\end{blocLivre}
\begin{blocLivre}{ex-1-15}
\begin{exercice}[3]{Une famille de paraboles}
Pour un réel $m$, on pose $P_m(x)=x^2-2mx+m+2$.
\partieExo{A}{Étude des racines}
\begin{questions}
\item Calculer le discriminant et le factoriser.
\item Déterminer les valeurs de $m$ pour lesquelles $P_m$ possède deux racines distinctes,
une racine double, ou aucune racine réelle.
\end{questions}
\partieExo{B}{Un point commun et un cas particulier}
\begin{questions}[resume]
\item Montrer que toutes les paraboles de la famille passent par un même point d'abscisse $\frac12$.
\item Pour $m=3$, résoudre $P_3(x)\leq0$ et déterminer le sommet.
\item Pour quelles valeurs de $m$ a-t-on $P_m(x)>0$ pour tout réel $x$ ?
\end{questions}
\end{exercice}
\end{blocLivre}
